\section{Limitations}
\label{sec:limitations}

\subsection{The Data Is American Sign Language}

The only large labelled dataset available to the project is the ASL Alphabet dataset. Every reported result in Deliverable~2 therefore describes American Sign Language, not Pakistan Sign Language. The two are different sign languages with different hand shapes, so the trained CNN cannot be used as a PSL recogniser and its scores do not transfer to PSL. The ASL baseline is used to build, train and compare the pipeline, not to make a claim about PSL accuracy.

\subsection{No PSL Photographs}

At the time of writing the team has collected 0 PSL photographs from 0 volunteers. The PSL letter set is also not final: published sources report 34 to 37 letters, and the current label map contains the 26 ASL letters A to Z. A PSL model, a PSL test set and any statement about PSL performance are therefore outside what this deliverable can show.

\subsection{A Subset of the Available Images}

Training does not use all 87{,}000 images of the ASL baseline. The shared split uses 400 cleaned images per class (10{,}400 in total, 7{,}280 in train, 280 per class in train) chosen by the data lead, and no per-class cap is applied at training time. This keeps tuning and final training feasible on the available hardware, but it reduces the variety the model sees and makes the results depend on which images were kept.

An existing PSL image dataset, UAlpha40 (CC BY 4.0, 36 static signs), is a candidate for PSL work in later deliverables; its folder structure and subject IDs have not been checked, and it was not used for any result in this report.

\subsection{Frames and Signers Are Not Independent}

The ASL baseline consists of video frames, and the number of distinct signers behind the dataset has not been verified. The shared split was generated by running the data lead's pipeline (\texttt{data\_scripts.build\_split\_pool} followed by \texttt{signbridge.data.splits}) on the ASL training folder with the committed configuration (seed 42, 400 images per class), giving 10{,}400 images in total (7{,}280 train, 1{,}560 validation and 1{,}560 test). Near-duplicates were removed from the split pool at a Hamming distance of 3 bits, the remaining images were placed in leader groups at 10 bits, and every group was assigned whole to one split, so the per-class shares are exactly 70/15/15. Grouping reduces leakage but does not remove it. Using the same 64-bit difference hash as the cleaning scripts, 21.0\% of validation and 20.3\% of test images have a training image of the same class within 4 bits; the figures rise to 49.9\% and 47.2\% within 6 bits, 73.4\% and 70.0\% within 8 bits, and 87.4\% and 87.4\% within 10 bits. The closest training neighbour is 4 bits away in both splits. The cause is that leader grouping bounds each image's distance to its own group leader, not the distance between every pair of images in different groups. The reported validation and test scores are therefore optimistic, and they are not an estimate of accuracy on new signers. One class also has a single group holding 43 of its 400 images (10.75\%), slightly above the 10\% target used when choosing the setting; all other classes are below 10\%.

\subsection{Small Compute Budget}

The long runs were done on a free cloud T4 GPU, not on the project laptop. The tuning search ran 12 trials of 3 epochs, about 26 minutes of training in total, and the final run stopped after 9 of the 12 configured epochs, about 6 minutes. The GPU took about 0.15 seconds per batch at $224\times224$ and batch size 32, against the 3.2 seconds per batch measured on the CPU, which is why the search and the final run were moved off the laptop. A larger budget would allow more trials, repeated seeds to measure run-to-run noise, more epochs or a larger subset. The split still uses 280 training images per class, so the results are limited by the data as well as by the search.

\subsection{Static Letters Only}

The system classifies a single image of one hand. It does not use motion, temporal information, two-handed signs, words or sentences, and it does not translate between PSL and text or speech. Letters that require movement are outside its scope.

\subsection{No Demographic Bias Assessment}

The dataset has no demographic labels, and no PSL volunteers were photographed. It is therefore not possible to measure whether the recogniser performs differently across age, gender, skin tone or signing background. The bias discussion in this report is limited to what can be measured without those labels, such as performance by letter and by image brightness or blur, and it cannot support a claim of fairness.
